\documentclass[12pt,a4paper]{article}

% ---------------------------------------------------------
% Packages
% ---------------------------------------------------------

\usepackage[utf8]{inputenc}
\usepackage[T1]{fontenc}
\usepackage{lmodern}
\usepackage{geometry}
\usepackage{graphicx}
\usepackage{float}
\usepackage{fancyhdr}
\usepackage{titlesec}
\usepackage{hyperref}
\usepackage{xcolor}
\usepackage{enumitem}
\usepackage{listings}

% ---------------------------------------------------------
% Page layout
% ---------------------------------------------------------

\geometry{
a4paper,
left=25mm,
right=25mm,
top=25mm,
bottom=25mm
}

% ---------------------------------------------------------
% Header and footer
% ---------------------------------------------------------

\pagestyle{fancy}
\fancyhf{}

\setlength{\headheight}{14.5pt}

\fancyfoot[L]{%
\footnotesize\ttfamily
\href{mailto:rostand.assonfo-demaboi@thga.stud.de}
{rostand.assonfo-demaboi@thga.stud.de}
}

\fancyfoot[C]{\footnotesize THGA Bochum}
\fancyfoot[R]{\footnotesize\thepage}

\renewcommand{\headrulewidth}{0pt}
\renewcommand{\footrulewidth}{0.4pt}

% ---------------------------------------------------------
% Code listings
% ---------------------------------------------------------

\lstset{
basicstyle=\ttfamily\small,
frame=single,
breaklines=true,
columns=fullflexible,
keepspaces=true
}

% ---------------------------------------------------------
% Hyperlinks
% ---------------------------------------------------------

\hypersetup{
colorlinks=true,
linkcolor=black,
urlcolor=blue,
pdftitle={User Manual -- Football Tournament Management System},
pdfauthor={Rostand Cedrix Assonfo Demaboi},
hypertexnames=false
}

% ---------------------------------------------------------
% Document
% ---------------------------------------------------------

\begin{document}

% =========================================================
% Title page
% =========================================================

\begin{titlepage}
\centering

\vspace*{2cm}

{\Huge\bfseries
User Manual\\[0.5cm]
Football Tournament Management System
}

\vspace{2cm}

{\Large
Rostand Cedrix Assonfo Demaboi
}

\vspace{0.5cm}

{\large
Summer Term 2026
}

\vfill

{\large
Technische Hochschule Georg Agricola, Bochum
}

\vspace{0.5cm}

{\large
Module: Introduction to Database Management Systems
}

\vspace{0.5cm}

{\large
Dozent: Stephan Bökelmann
}

\vfill

{\large
\today
}

\end{titlepage}

% =========================================================
% Table of contents
% =========================================================

\pagenumbering{roman}

\clearpage
\tableofcontents
\clearpage

\pagenumbering{arabic}

% =========================================================
% 1 Introduction
% =========================================================

\section{Introduction and System Overview}

The Football Tournament Management System is an application
designed to organize, track, and administer football
competitions.

The system follows a decoupled architecture consisting of
three main components:

\begin{itemize}
\item \textbf{Backend:} A REST API developed with FastAPI
and connected to a relational database.

\item \textbf{Frontend:} A desktop graphical interface
developed with Python Tkinter for user interaction.

\item \textbf{Orchestration:} Docker Compose is used to
simplify deployment and ensure a reproducible runtime
environment.
\end{itemize}

The application provides two main access modes.

\subsection{Organizer Mode}

Organizer Mode provides access to administrative operations.

Authenticated organizers can create, modify, and delete
tournaments, teams, and matches.

Authentication is performed using an API token.

\subsection{Viewer Mode}

Viewer Mode provides read-only access to the application.

Users can view tournaments, teams, matches, scores, and
rankings without modifying the database.

% =========================================================
% 2 Backend Deployment
% =========================================================

\section{Backend Deployment and Launch}

The backend server and database can be started using Docker
Compose or the backend can be launched separately for
development and testing.

This section describes both methods.

\subsection{Prerequisites}

Before starting the application, make sure that the following
software is installed:

\begin{itemize}
\item Docker
\item Docker Compose
\item Python 3
\end{itemize}

\subsection{Deployment Using Docker Compose}

Open a terminal and navigate to the project root directory:

\begin{lstlisting}
cd ~/football-tournament-manager
\end{lstlisting}

Start the backend and database containers:

\begin{lstlisting}
docker compose up --build -d
\end{lstlisting}

The \texttt{--build} option ensures that the required Docker
images are rebuilt when necessary, while \texttt{-d} starts
the services in detached mode.

Check the status of the containers with:

\begin{lstlisting}
docker compose ps
\end{lstlisting}

The containers should be displayed as running.

The FastAPI backend is available on port:

\begin{lstlisting}
8000
\end{lstlisting}

The default local API URL is:

\begin{lstlisting}
http://localhost:8000
\end{lstlisting}

\subsection{Separate Backend and Frontend Launch}

The application can also be started without using Docker
Compose for the frontend.

In this configuration, the backend is started in one terminal
and the frontend is started in a second terminal.

This method is useful for development and testing because the
frontend and backend can be started and monitored independently.

\subsubsection{Terminal 1 -- Start the Backend}

Open the first terminal and navigate to the backend directory:

\begin{lstlisting}
cd ~/football-tournament-manager/backend
\end{lstlisting}

Activate the Python virtual environment if one is configured
for the backend:

\begin{lstlisting}
source .venv/bin/activate
\end{lstlisting}

Start the FastAPI backend with:

\begin{lstlisting}
uvicorn app.main:app --reload --host 0.0.0.0 --port 8000
\end{lstlisting}

Keep this terminal open while using the application.

The backend should then be available at:

\begin{lstlisting}
http://localhost:8000
\end{lstlisting}

\subsubsection{Terminal 2 -- Start the Frontend}

Open a second terminal.

Navigate to the frontend directory:

\begin{lstlisting}
cd ~/football-tournament-manager/frontend
\end{lstlisting}

Activate the frontend virtual environment:

\begin{lstlisting}
source .venv/bin/activate
\end{lstlisting}

Start the graphical application:

\begin{lstlisting}
python3 src/football_tournament/main.py
\end{lstlisting}

The frontend connects to the backend running in the first
terminal.

In the connection dialog, use:

\begin{lstlisting}
http://localhost:8000
\end{lstlisting}

For Organizer Mode, enter:

\begin{lstlisting}
mykey
\end{lstlisting}

For Viewer Mode, leave the API token empty or enter an
incorrect token.

\subsection{Difference Between the Two Launch Methods}

There are two possible ways to run the application.

With Docker Compose:

\begin{lstlisting}
docker compose up --build -d
\end{lstlisting}

Docker manages the backend and database services.

With separate terminals, the backend and frontend are started
independently.

The separate method is particularly useful during development
because changes to the Python backend or frontend can be tested
directly.

It is important to ensure that both launch methods use the same
database configuration if identical initial data is expected.

% =========================================================
% 3 Frontend
% =========================================================

\section{Frontend Launch and Connection}

Once the backend is running, the desktop frontend can be
started.

\subsection{Activating the Python Environment}

Navigate to the frontend directory:

\begin{lstlisting}
cd ~/football-tournament-manager/frontend
\end{lstlisting}

The project uses a virtual environment named
\texttt{.venv}. Activate it with:

\begin{lstlisting}
source .venv/bin/activate
\end{lstlisting}

Then start the application with:

\begin{lstlisting}
python3 src/football_tournament/main.py
\end{lstlisting}

\subsection{Startup Connection Dialog}

When the application starts, a connection dialog is displayed.

The dialog requests the following information:

\begin{itemize}
\item \textbf{API URL:} The address of the FastAPI backend.
The default value is
\texttt{http://localhost:8000}.

\item \textbf{X-API-Token:} The authentication token used
to access organizer operations.
\end{itemize}

\subsection{Viewer Mode}

To access the application in read-only mode, leave the
\texttt{X-API-Token} field empty or enter an incorrect API token.

In both cases, the application starts in Viewer Mode.

\subsection{Organizer Mode}

To access administrative functions, enter the API token
configured for the application.

The API token used by the application is:

\begin{lstlisting}
mykey
\end{lstlisting}

Only the correct token allows access to Organizer Mode.

If the token is empty or incorrect, the application connects
in Viewer Mode.

% =========================================================
% 4 Interface
% =========================================================

\section{Interface Navigation and Usage}

The main interface is organized as a tabbed dashboard.

The available sections allow users to inspect and manage the
main entities of the football tournament system.

\subsection{Tournaments Tab}

The Tournaments tab displays the registered football
competitions.

Users can view information such as:

\begin{itemize}
\item Tournament name
\item Start date
\item End date
\item Current status
\end{itemize}

Organizers can create, modify, or delete tournaments when
authenticated with a valid API token.

\subsection{Teams Tab}

The Teams tab displays the participating teams.

For each team, the interface provides information such as:

\begin{itemize}
\item Unique team identifier
\item Team name
\item Coach or manager name
\end{itemize}

Organizers can create, update, and delete teams.

The current application interface does not provide a dedicated
player-management workflow in this section. Therefore, players
should not be considered a directly managed feature of the
current graphical interface.

\subsection{Matches Tab}

The Matches tab displays the fixtures registered in the system.

A tournament selector at the top of the tab allows users to
filter the displayed matches.

The available options include:

\begin{itemize}
\item A specific tournament
\item All tournaments
\end{itemize}

The match table displays information such as:

\begin{itemize}
\item Match date
\item Home team
\item Away team
\item Score
\item Referee
\end{itemize}

For a match that has not yet received a score, the score is
displayed using a dash (\texttt{-}).

Validated scores are displayed using the
Home--Away format.

\subsection{Standings Tab}

The Standings tab displays the current ranking of teams in a
tournament.

The ranking is calculated dynamically from the recorded match
results.

The displayed statistics include:

\begin{itemize}
\item Points
\item Wins
\item Draws
\item Losses
\item Goals scored
\item Goals conceded
\item Goal difference
\end{itemize}

The standings therefore reflect the current results stored in
the database.

% =========================================================
% 5 Administrative Operations
% =========================================================

\section{Administrative Operations}

Administrative operations are available to users authenticated
in Organizer Mode.

The following operations can be performed from the graphical
interface.

\subsection{Scheduling a Match}

To create a new match:

\begin{enumerate}
\item Open the \textbf{Matches} tab.
\item Click the \textbf{+ Schedule Match} button.
\item Select the tournament.
\item Select the home team.
\item Select the away team.
\item Select or enter the match date.
\item Enter the referee information if required.
\item Confirm the operation using the
\textbf{SAVE MATCH} button.
\end{enumerate}

The new match is then stored in the database and becomes
available in the Matches tab.

\subsection{Updating a Match Score}

To update the score of an existing match:

\begin{enumerate}
\item Open the \textbf{Matches} tab.
\item Select the desired match.
\item Click \textbf{Update Score}.
\item Enter the goals scored by the home team.
\item Enter the goals scored by the away team.
\item Confirm the modification.
\end{enumerate}

The new result is stored in the database.

The standings are subsequently recalculated from the updated
match results.

\subsection{Deleting Data}

Organizers can delete obsolete teams or matches using the
corresponding \textbf{Delete} operation.

Deletion is subject to the database's referential integrity
constraints.

If an object is referenced by another database object, the
database may prevent the deletion in order to preserve data
consistency.

\subsection{Refreshing the Data}

The application provides a \textbf{Refresh Data} function.

This operation retrieves the current data from the backend API
and updates the local graphical display.

It is particularly useful when data has been modified externally
or by another client.

\subsection{Logout and Switching User Mode}

To terminate the current session, close the application or use
the available disconnect/logout functionality in the connection
panel.

When the application is started again, the connection dialog is
displayed and the user can choose whether to connect in Viewer
Mode or Organizer Mode.

% =========================================================
% 6 Typical Usage Workflow
% =========================================================

\section{Typical Usage Workflow}

A typical tournament-management workflow can be performed in the
following order:

\begin{enumerate}
\item Start the backend and database using Docker Compose,
or start the backend separately in the first terminal.
\item Verify that the backend is running.
\item Open a second terminal and start the frontend
application.
\item Enter the backend API URL.
\item Enter the API token to access Organizer Mode, or leave
it empty for Viewer Mode.
\item Create a tournament.
\item Create the participating teams.
\item Schedule the matches between the teams.
\item Enter the results after matches have been played.
\item Open the Standings tab to view the updated ranking.
\item Use Refresh Data whenever information needs to be
synchronized with the backend.
\end{enumerate}

This workflow allows an organizer to manage a complete football
competition from its creation through the recording of match
results and rankings.

% =========================================================
% 7 Troubleshooting
% =========================================================

\section{Troubleshooting}

\subsection{Docker Containers Are Not Running}

If the containers are not running, execute:

\begin{lstlisting}
docker compose ps
\end{lstlisting}

To inspect the container logs, use:

\begin{lstlisting}
docker compose logs
\end{lstlisting}

The backend can also be restarted with:

\begin{lstlisting}
docker compose up --build -d
\end{lstlisting}

\subsection{Frontend Cannot Connect to the Backend}

First verify that the backend is running.

The default API address is:

\begin{lstlisting}
http://localhost:8000
\end{lstlisting}

Also verify that the API token is correct when using Organizer
Mode.

\subsection{Backend Started Separately}

If the backend is started manually, open the first terminal and
execute:

\begin{lstlisting}
cd ~/football-tournament-manager/backend
\end{lstlisting}

Then start the FastAPI application:

\begin{lstlisting}
uvicorn app.main:app --reload --host 0.0.0.0 --port 8000
\end{lstlisting}

Keep this terminal running.

Open a second terminal and execute:

\begin{lstlisting}
cd ~/football-tournament-manager/frontend
source .venv/bin/activate
python3 src/football_tournament/main.py
\end{lstlisting}

The frontend must connect to:

\begin{lstlisting}
http://localhost:8000
\end{lstlisting}

\subsection{Virtual Environment Cannot Be Activated}

Make sure that you are inside the frontend directory:

\begin{lstlisting}
cd ~/football-tournament-manager/frontend
\end{lstlisting}

Then activate the correct environment:

\begin{lstlisting}
source .venv/bin/activate
\end{lstlisting}

If the virtual environment does not exist, it must be created
according to the project's Python environment setup instructions.

% =========================================================
% 8 Conclusion
% =========================================================

\section{Conclusion}

The Football Tournament Management System provides a complete
interface for managing football tournaments, teams, matches,
scores, and standings.

The system separates the graphical frontend from the REST API
and database layer. Docker Compose provides a reproducible
environment for running the backend and database.

The application can also be launched separately, with the
backend running in one terminal and the frontend running in a
second terminal.

Viewer Mode allows participants to inspect tournament
information, while Organizer Mode provides the administrative
operations required to maintain the competition.

For development and deployment, the most important components
are the Docker Compose environment, the FastAPI backend, the
Python virtual environment, and the Tkinter frontend.

\end{document}